\section{Step 25 Round a diagonal prediction once}\label{h263-prediction-s05}
For two active phases, load all four neighbours and form their sum before adding
two and shifting by two. Retain the unrounded intermediate. Do not reuse a helper
that has already rounded a horizontal average unless it preserves the information
needed for the final exact result.

Check the four-neighbour counterexample from the figure, then a high-valued
pattern. A sum of four 255s is 1020 and still predicts 255 after the bias and shift.
Both examples distinguish the width of a sample from the width of an intermediate.

\paragraph{Observe} With three zeros and one one, diagonal prediction is zero.
Two rounds of upward pair averaging would incorrectly produce one.
\Needspace{8\baselineskip}
\paragraph{Check} When is the diagonal rounding bias applied?
\begin{itemize}
  \item[A.] Once, after summing all four neighbours.
  \item[B.] Separately after every addition.
  \item[C.] Only after the signed residual has been clipped to a byte.
\end{itemize}
